% ==========================================================================
% Paper 2 — Section 2: Formal Framework. Five Primitives, the Frozen
% Baseline, the Extension-Operator Catalog, and the Audit Methodology.
% ==========================================================================
\section{Formal Framework: Five Primitives and the Frozen Baseline}
\label{sec:framework}

This section fixes the definitions that the remainder of the paper uses
without further negotiation: the five computational primitives (§2.2), the
frozen dynamical baseline they are tested against (§2.3), the catalog of
single-degree-of-freedom extensions that ascend the ladder (§2.4), and the
audit methodology by which every rung is either established or rejected
(§2.5). All quantitative constants in this section are taken verbatim from
the frozen archives (TAN-I and TAN-II sprints 4.1--4.3-A); none are free.

\subsection{Notation and conventions}
\label{sec:notation}

Time is discrete, $t\in\mathbb{N}$. Inputs are scalar amplitudes
$x_t\in\mathbb{R}$; the temporal window is the ordered $W$-tuple
$X_t=(x_{t-W+1},\ldots,x_t)$ with zero padding at onset. We write
$[z]_+=\max(z,0)$ for rectification, $\Theta(\cdot)$ for the Heaviside
step, $\sigma(z)=1/(1+e^{-z})$ for the logistic function, and
$\operatorname{softmax}_i(\bm e)$ for the normalized exponential map.
In Sprint~4.2-B/C/D the standard max-stabilized form
$\alpha_i=\exp(e_i-\max_k e_k)/\sum_j\exp(e_j-\max_k e_k)$ is evaluated
directly without offset ($\sum_i\alpha_i=1.0$ exactly), whereas the baseline
model implementation (\texttt{tan.py}) added a numerical stabilizer
$\delta=10^{-9}$ to the denominator. Alternative mappings that induce exact
sparsity without infinite gain, such as sparsemax \citep{martins2016}, are
discussed as theoretical comparisons. Candidate sets are always
finite. A \emph{counterfactual pair} is a pair of experimental conditions
that share the identical history content and differ only in the query; a
\emph{negative-control chain} is a sequence of models $M_0,M_1,M_2$ in which
each model adds exactly one organizational degree of freedom and the
preceding model is required to fail the criterion that the following model
is required to pass.

\subsection{The five primitives, operationally defined}
\label{sec:primitives}

The paper's object is the chain
$\text{Memory}\to\text{Saliency}\to\text{Routing}\to\text{Binding}\to
\text{Composition}$. Each primitive is defined by an operational criterion
that a frozen experiment can pass or fail; the criteria are the ones
actually used in the audits of §4--§8. The chain is constructive: in the
mechanism family tested here, each later primitive presupposes the earlier
ones, while the central negative results of this paper show that none of
the reverse implications holds.

\begin{description}[style=nextline,leftmargin=2em]
\item[P1 --- Memory.] \emph{Definition.} A system has memory relative to a candidate
state variable if its future evolution depends causally on aspects of the input history
that are not recoverable from that candidate state. \emph{Criterion} (TAN-I Phase 1, §3):
the one-dimensional scalar membrane potential $h_t$ is not state-sufficient; the projection
$(h_t,X_t)\mapsto h_t$ is not lumpable. Two histories with equal membrane potentials but
different window contexts produce different downstream trajectories (quantified by pairwise
separation $K$ and post-reset divergence). This demonstrates state non-sufficiency of scalar $h_t$;
it does not imply that the complete Markovian state $(h_t, X_t)$ violates Markovianity, nor
that Markovian dynamics imply state-space contraction.
\item[P2 --- Saliency weighting.] \emph{Definition.} A system performs saliency weighting
if its readout over a fixed candidate set preserves relative ranking across all active query
levels. Under standard continuous softmax (whose support is everywhere strictly positive,
in contrast to sparse activation functions such as sparsemax \citep{martins2016}), the defining
signature of scalar modulation is not support-set invariance, but rather \emph{ranking invariance}:
for all $S>0$, the induced ordering of attention weights $\alpha(Q)$ is strictly invariant,
and $\operatorname{argmax}_i\alpha_i(Q)$ is independent of query content. \emph{Criterion:}
Rank-1 kernel factorizability $E(Q, K) = q(Q)\cdot k(K)$ is a sufficient condition ensuring
that positive scalar query modulation acts purely as a monotonic gain on logit differences,
\item[P3 --- Content routing.] \emph{Definition.} A system routes if there
exist queries $Q_A,Q_B$ and a fixed history such that
$\operatorname{argmax}_i R(Q_A,K_i)\neq\operatorname{argmax}_i R(Q_B,K_i)$.
In multi-candidate memory ($N\ge3$), true content addressing requires selective
identifiability: for every candidate key $K_t$, there exists a query making $K_t$ the
unique maximum. While separable dot-product kernels $E(Q, K) = Q K$ in $\mathbb{R}$
can at most invert global order between the extreme elements without addressing intermediate keys
(Proposition~\ref{prop:scalar_impossible}), metric compatibility kernels $E(Q, K) = -(Q-K)^2$
achieve universal addressability even in $d=1$ (Proposition~\ref{prop:metric_suff}).
\emph{Criterion:} the counterfactual flip rate
$\mathrm{FlipRate}_{\mathrm{CF}}$ is nonzero on a fixed-history ensemble;
capacity (the geometric existence of flips) and realization (the running
model actually flipping) are reported separately (\S\ref{sec:routing}).
\item[P4 --- Semantic binding.] \emph{Definition.} A system binds if
role/address (key) and content entity (value) are decoupled and the output
transmits the value of the \emph{selected} key, following the
(key,value) association rather than key identity, position, or amplitude.
\emph{Criterion:} the value-swap counterfactual is antisymmetric,
$T_{\mathrm{swap}}=-T$ exactly, and the softness distance
$D=|C-v_{\mathrm{winner}}|$ is measured; hard binding is the limit $D=0$.
\item[P5 --- Composition.] \emph{Definition.} A system composes if it applies
a deterministic algebraic relation $f$ to at least two independently
retrieved bound values, $y=f(V_A,V_B)$. \emph{Criterion:} the three-error
joint audit of §2.5 (binding errors and composition error simultaneously
$\approx0$) together with the counterfactual suite of §8 (query swap, value
swap, key swap, channel-output swap, the zero-value trap, and the
binding-preserving, binding-breaking, and relational controls); conditioned
on correct binding the composition error must be exactly zero (the
binding--combiner separation test).
\end{description}

\subsection{The frozen baseline: the scalar Temporal Attention Neuron}
\label{sec:baseline}

The substrate is the TAN-I Temporal Attention Neuron, frozen before this
study and never retrained or retuned here. Its computation chain is

\begin{align}
\mu_t &= \frac{1}{W}\sum_{k=0}^{W-1} x_{t-k}, \qquad
S_t=[x_t-\mu_t-\varepsilon]_+,
\label{eq:surprise}\\
E_{t,i}&=\beta\,W_q W_k\,S_t\,x_{t-i}, \qquad
\alpha_{t,i}=\frac{\exp(E_{t,i})}{\sum_j\exp(E_{t,j})+\delta},
\label{eq:kernel}\\
C_t&=\sum_i \alpha_{t,i}\,x_{t-i}, \qquad
A_t=\tanh(S_t)\,C_t,
\label{eq:context}\\
h_t&=\lambda\,h_{t-1}+A_t, \qquad
y_t=\Theta(h_t-\theta_{\mathrm{sp}}), \qquad
h_t\leftarrow 0\ \text{on spike}.
\label{eq:membrane}
\end{align}

Three properties of this chain carry the whole study and must be kept in
view. (i) The query of the kernel \eqref{eq:kernel} is the non-negative
scalar surprise $S_t$ (rectified by \eqref{eq:surprise} and multiplied by
positive constants), so the query domain is the closed cone
$\mathcal{Q}_+=[0,\infty)$; the kernel is therefore of the separable form
(query scalar)$\times$(key scalar), which pins the attention argmax to the
key magnitudes regardless of the query---this is the exact content of the
order-conservation theorem (§5) and the reason TAN-I froze at
$\mathrm{FlipRate}=0.000$. (ii) The readout \eqref{eq:context} is a convex
combination, so the delivered ``value'' is always a mixture; discreteness is
not available to this substrate until a selection operator is added
(§7). (iii) The membrane update \eqref{eq:membrane} contains the discrete
spike-and-reset map; the system is piecewise discontinuous and must not be
treated as a smooth autonomous dynamical system (§4's conditional phase
planes are labeled accordingly).

Table~\ref{tab:params} is the frozen parameter inventory. Provenance
columns refer to the frozen archive; the TAN-II ladder re-freezes the TAN-I
values wherever it does not add a new degree of freedom, and every deviation
is pre-registered (for example, the default spike threshold of the E-I
stages is $\theta_{\mathrm{sp}}=3.0$, chosen pre-run to keep the tuning
criterion decisive, with the TAN-I value $0.5$ retained in the threshold
sweep).

\begin{table}[t]
\centering
\small
\begin{tabular}{@{}lp{4.9cm}p{2.6cm}l@{}}
\toprule
Symbol & Frozen value & Meaning & Provenance \\
\midrule
$W$ & $5$ & temporal window & TAN-I \\
$\lambda$ & $0.5$ & membrane leak & TAN-I \\
$\theta_{\mathrm{sp}}$ & $0.5$ (TAN-I); $3.0$ default in TAN-II (sweep
$\{0.5,1.0,3.0,6.0\}$) & spike threshold & TAN-I / 4.1 \\
$\varepsilon,\varepsilon_I$ & $0,\,0.3$ & surprise dead zones & TAN-I / 4.1 \\
$\beta, W_q, W_k$ & $1.0,\,2.0,\,1.0$ & kernel constants
($c=\beta W_q W_k=2.0$) & TAN-I \\
$\delta$ & $10^{-9}$ & softmax offset & TAN-I \\
$(w_{EI},w_{IE},w_{II})$ & $(0.8,1.0,0.2)$ & Dale-signed couplings & 4.1 \\
$(w_E,w_I,\theta_E,\theta_I)$ & $(0.5,1.0,0.4,1.0)$ & static opponent
thresholds (C1) & 4.1 \\
$\omega$ & $0.4$ & query-direction rotation & 4.2-B \\
$(Q_A,Q_B)$ & $(3.0,4.0)$ & counterfactual queries (single-query) & TAN-I/4.2 \\
$(Q_1,Q_2)$ & $(4.5,5.3)$ & dual-query window & 4.3-A \\
$(V_A,V_B)$ scenarios & $(5,9)$, $(7,-2)$, $(2,6)$ & canonical values & 4.2-C/4.3-A \\
$\gamma$ ladder & $\{1,2,5,10,100,\infty\}$ & WTA sharpening & 4.2-D \\
Seeds / $N$ / grid & $20260904$--$06$; $10^4$; $2001^2$ & Monte Carlo /
quadrature & 4.2--4.3 \\
\bottomrule
\end{tabular}
\caption{Frozen parameter inventory. All values are pre-registered; the
provenance column names the frozen sprint. The spike threshold column shows
the two frozen conventions and the sweep that audits their consequences.}
\label{tab:params}
\end{table}

\begin{figure}[t]
\centering
\includegraphics[width=\textwidth]{fig1_taxonomy}
\caption{The conceptual taxonomy and the frozen TAN architecture.
(a) The five primitives of §2.2 as a constructive chain. (b) The frozen
TAN-I computation chain (§2.3). (c) The degenerate geometry of the scalar
kernel: the non-negative query cone $\mathcal{Q}_+$ leaves the candidate
ordering locked to the key magnitudes (Theorem 5.1).}
\label{fig:taxonomy}
\end{figure}

\subsection{The catalog of extension operators}
\label{sec:operators}

Each rung of the ladder is one of the following operators, applied to the
baseline and left unchanged afterwards. No rung changes anything except the
named operator.

\begin{description}[style=nextline,leftmargin=2em]
\item[O1 --- Opponent competition (Sprint 4.1).] Two cells with membrane
potentials
$h_E(t)=\lambda_E h_E(t-1)+A_E(t)-w_{EI}A_I(t)$ and
$h_I(t)=\lambda_I h_I(t-1)+w_{IE}A_E(t)-w_{II}A_I(t)$, where $A_E,A_I$ are
the per-branch drives (attention-gated contexts in the attended models;
static thresholded currents $[w_E x_t-\theta_E]_+$, $[w_I x_t-\theta_I]_+$
in the static control C1) and all couplings obey Dale's sign structure
$w_{ab}\ge0$. The four-model ablation ladder C1$\to$C2$\to$C3$\to$C4 of
§4 instantiates this operator.
\item[O2 --- Vector query--key interaction (Sprint 4.2-B).] Keys are
embedded as unit vectors $\hat\varphi_d(x)=p_d(x)/\|p_d(x)\|$ with
$p_d(x)=(x,x^2,\ldots,x^d)$ and $\hat\varphi_d(0)=0$; the query is the
rotating vector $Q_t=c\,S_t\,\hat u(\theta_t)$ with angle $\theta_t=\omega S_t$ and
$\hat u(\theta)=(\cos\theta,\sin\theta,0,\ldots)$; the kernel becomes
$E_{t,i}=Q_t^{\top}K_i$. Only the representation changes: window, surprise
gate, dynamics, and readouts are frozen.
\item[O3 --- Key--value decoupling (Sprint 4.2-C).] Each slot carries a
pair $(k_i,v_i)$: keys drive the attention, values are transmitted,
$C_t=\sum_i\alpha_{t,i}v_i$. Query slots carry $v=0$ and are never
candidates.
\item[O4 --- Readout discreteness (Sprint 4.2-D).] The selection function
$\alpha_i(\gamma)=\operatorname{softmax}_i(\gamma\cdot\mathrm{logit}_i)$
with $\gamma\in[1,\infty]$; $\gamma=\infty$ is the one-hot
winner-take-all.
\item[O5 --- Parallel channels and the combiner (Sprint 4.3-A).] The
frozen window is replaced by the dual-query window $[A,0,B,Q_1,Q_2]$ with
$(Q_1,Q_2)=(4.5,5.3)$ (the direction-targeting derivation is frozen in the
pre-registration); channel $c\in\{1,2\}$ is keyed to slot $3+c$ with the
shared window mean, and a two-input node computes $y=f(\hat C_1,\hat C_2)$
with $f\in\{+,-\}$.
\end{description}

\subsection{Audit methodology and discipline}
\label{sec:audit}

Every rung is established or rejected by the following machinery, which is
itself part of the frozen protocol and is reported in §10 together with the
claim ledger.

\emph{Counterfactual suites.} A result is claimed only if it survives the
counterfactual manipulations of its rung: query swap, value swap, key swap,
channel-output swap, position randomization, amplitude randomization, the
single-query controls, and the leakage checks ($Q\cap V=\varnothing$;
queries are never candidates; equal-norm keys). Subtraction provides the
non-commutative operator whose swap responses are antisymmetric, addition
the commutative one whose swap responses are invariant---both predictions
are tested, so that correct commutative behavior can never be scored as a
false negative.

\emph{The three-error joint audit.} Composition success requires
$E_{\mathrm{bind},A}\approx0$ \emph{and} $E_{\mathrm{bind},B}\approx0$
\emph{and} $E_{\mathrm{comp}}=|y-f(V_A,V_B)|\approx0$ simultaneously. The
joint requirement exists to exclude the accidental-zero class of false
positives exemplified by the $V_A=0$ trap of §8, in which the composition
error vanishes while a binding error remains.

\emph{Negative-control chains and orthogonal axes.} Each positive claim is
paired with the demonstration that the previous rung fails the same
criterion ($M_0$ fails composition, $M_1$ binds but does not compose,
$M_2$ composes). Two orthogonal axes are formalized from these chains:
\emph{routing selectivity $\perp$ readout discreteness} (§7: the WTA ladder
changes discreteness while the flip rate stays exactly flat) and
\emph{binding fidelity $\perp$ combiner error} (§8: the combiner adds zero
error conditioned on correct binding).

\emph{Statuses and stop rules.} Every check receives an explicit
\PassFailInconclusive{}. Discrepancies with frozen predictions trigger the
stop-and-debug procedure: stop, independently recompute, record the cause in
a dated amendment that preserves the original prediction verbatim, and only
then continue. The six amendments of this programme (AM1--AM6: the E4(iii)
closed-form value, the unnormalized-control expectation, the corrected
fixed-point constants, the E-I drive-transfer value, the readout-softness
monotonicity, and the canonical routing margins) are all recorded this way;
none altered a model, query, seed, or threshold.

\emph{Determinism and statistical criteria.} All experiments are
deterministic under the fixed seeds $20260904$--$20260906$ ($N=10^4$ per
seed); analytic benchmarks are deterministic $2001^2$ quadratures computed
from the frozen model. Probability estimates are accepted when
$|\hat p-p_{\mathrm{bench}}|\le3\,\mathrm{SE}+0.002$ and mean-value
estimates when $|\hat m-m_{\mathrm{bench}}|\le3\,\mathrm{SE}_{\mathrm{mean}}
+0.01$; structural identities (swap antisymmetry, position invariance,
conditioned composition error) are asserted at machine precision
($10^{-12}$).
